\chapter{Síntesis integral de COBRIA}

COBRIA organiza decisiones que aparecen cuando un sistema NoSQL conserva documentos
operacionales y necesita búsquedas, resúmenes, conjuntos analíticos y herramientas de
inteligencia artificial. Existe una fuente canónica por hecho de dominio y toda
representación derivada declara cómo se obtiene, verifica, delimita y reconstruye.

\section{Respuesta a la investigación}

La estructura Objeto--Dater--Servicio--Catálogo no constituye por sí sola una teoría
nueva. COBRIA aporta una composición explícita para documentos NoSQL: identidad y
revisión en el canónico; transformación y procedencia en el Dater; reglas y autorización
en servicios; semántica versionada en catálogos; y reconstrucción en cada proyección.

La reducción de lectura depende de consulta, índice, forma documental y distribución.
Una proyección se adopta después de compararla con la consulta canónica indexada. La
amplificación se registra mediante operaciones, bytes, eventos, reintentos,
almacenamiento y trabajo operacional.

La adaptabilidad hacia minería de datos e inteligencia artificial procede del linaje:
conjuntos, características, fragmentos e inferencias conservan fuente, tiempo, versión y
ámbito. Esto permite reconstruir un corte, detectar deriva y evitar que recuperación o
agentes crucen fronteras de autorización.

\chapter{Guía de adopción}

Considérese COBRIA cuando existen múltiples consumidores, documentos heterogéneos,
aislamiento organizacional, evolución frecuente, derivados costosos de reparar o
productos de datos reproducibles. Simplifíquese cuando una colección indexada satisface
el objetivo o el costo de gobierno supera el riesgo.

\begin{enumerate}[leftmargin=*]
  \item identificar la escritura autoritativa y cerrar rutas paralelas;
  \item hacer obligatorio el ámbito en acceso, caché, eventos y recuperación;
  \item añadir revisión, tiempo y versión de esquema;
  \item medir la consulta antes de crear una proyección;
  \item declarar transformación, equivalencia y reconstrucción;
  \item publicar por versión y practicar reparación;
  \item extender procedencia hacia analítica e inteligencia artificial;
  \item retirar derivados sin beneficio o responsable.
\end{enumerate}

\chapter{Contribuciones y límites}

La contribución conceptual es el dato reconstruible como contrato. La formal es la tupla
que une fuente, transformación, ámbito, versión, equivalencia y reparación. La
arquitectónica es una frontera coherente desde documentos hasta productos de IA. La
práctica es un catálogo, listas de comprobación y protocolo reproducible.

La contribución empírica conserva dos estratos. El primero contiene 180 corridas locales
históricas sobre PouchDB y LokiJS. El segundo contiene 270 corridas distribuidas
históricas sobre MongoDB 8.0, CouchDB 3.4 y OpenSearch 2.19.3. Todas superaron P1, P2,
R1, S1 y A1 bajo sus oráculos originales y no registraron fugas de ámbito. Como una
auditoría posterior fortaleció el arnés, estas cifras describen evidencia histórica y no
certificación vigente. La repetición preliminar corregida completó casos pequeños en
MongoDB y CouchDB; OpenSearch 3.2.0 quedó pendiente por incompatibilidad de plataforma.

Falta repetir las celdas distribuidas completas con el arnés corregido; ejecutar
OpenSearch 3.2.0 en ARM nativo o infraestructura compatible; realizar una réplica
independiente; completar el doble cribado bibliográfico; y ejecutar el estudio humano de
transferencia. P1, P2, R1, S1 y A1 forman el conjunto de criterios de conformidad, pero
la certificación independiente permanece pendiente.

\chapter{Conclusión}

COBRIA hace visibles seis preguntas: qué documento manda, a qué ámbito pertenece, qué
versión representa, qué copia puede perderse, cómo se reconstruye y cuánto cuesta
conservarla. Al extenderlas a conjuntos de datos, recuperación e inferencias, ofrece un
puente entre ingeniería NoSQL, minería de datos e inteligencia artificial.

La propuesta es suficientemente concreta para implementarse y refutarse. La réplica
embebida apoya sus contratos nucleares, mientras la evaluación administrada y humana
sigue abierta. Esa frontera define con honestidad qué se sabe y qué debe venir después.
